% !PRINTED 299
\weberpage{299}
\begin{center}
\textbf{\S\ 81. \gesperrt{Non-normal divisors of $\mathfrak{L}_0$.}}
\end{center}

We now inquire, leaving aside the case $p = 3$, after the non-normal divisors of the group $\mathfrak{L}_0$ whose index is smaller than $p+1$; upon the existence of such divisors depends the possibility of forming resolvents of the transformation equation \gesperrt{whose degree is lower than $p+1$.}

First of all it can be shown that the index of a divisor $\mathfrak{K}$ of $\mathfrak{L}_0$ can never be smaller than $p$.

Let
\weq{1}{\mathfrak{K} = S_0,\ S_1,\ \ldots,\ S_{\nu-1}}
be any divisor whatever of $\mathfrak{L}_0$ of degree $\nu$, and, if it is possible, let
\weq{2}{T = \begin{pmatrix} \alpha, & \beta \\ \gamma, & \delta \end{pmatrix}}
be a substitution of degree $p$ which occurs in $\mathfrak{L}_0$ but not in $\mathfrak{K}$. Then too, since $p$ is a prime, no power of $T$ lower than the $p$th can occur in $\mathfrak{K}$. The subsidiary groups
\weq{3}{\mathfrak{K},\ T\mathfrak{K},\ T^2\mathfrak{K},\ \ldots,\ T^{p-1}\mathfrak{K}}
contain none but \gesperrt{mutually distinct elements} $T^i S_k$, and accordingly $\nu p$ is at most equal to the degree of the group $\mathfrak{L}_0$, that is:
\[
\nu p \leqq p\,\frac{p^2-1}{2}.
\]
Thus the index of $\mathfrak{K}$, that is the quotient
\[
p\,\frac{p^2-1}{2} : \nu
\]
\gesperrt{is equal to $p$ or greater than $p$.}

A divisor $\mathfrak{K}$ of $\mathfrak{L}_0$ whose index is smaller than $p$ must therefore contain all substitutions of degree $p$, and hence also (§ 79, I.) all substitutions $T$ in which
\weq{4}{\alpha + \delta \equiv 2 \pmod{p}}
holds. Accordingly such a group $\mathfrak{K}$ contains in the first place the substitution
\weq{5}{A = \begin{pmatrix} 1, & 0 \\ 1, & 1 \end{pmatrix}}
%CONT>>
